\chapter{Discussion and Future Perspectives}
\label{ch:discussion}

This chapter interprets the results of \cref{ch:method}. It compares them with the original paper and
the wider literature (\cref{sec:disc-compare}), discusses their theoretical and practical implications
(\cref{sec:disc-implications}), states the limitations of the work (\cref{sec:disc-limits}) and proposes
directions for future research (\cref{sec:disc-future}).

\section{Comparison with previous work}
\label{sec:disc-compare}

\paragraph{Confirmations.} The label semantics of Zou et al.\ can be recovered from the public
annotations. That the reconstructed label distribution matches the paper's Fig.~3 for almost every
class suggests that the specificity hierarchy, left inheriting straight and right governed only by
dedicated signals were rebuilt correctly. A model trained on these labels reaches a TSS of 0.899 on
subset A val, in the range the paper reports for its direction-aware model (0.880). This confirms that
maneuver-level prediction from a single image is feasible without HD maps, which map-based
systems~\cite{fairfield2011traffic,possatti2019traffic,polley2025atlas} rely on, and without the
annotated lane metadata of earlier map-free work~\cite{langenberg2019deep}.

\paragraph{Differences.} The rebuilt benchmark contains about 29\% fewer frames than the paper's, for
reasons that cannot be determined without the original filtering rules. The paper describes its data as
Subset-A, which its own frame counts contradict; the evidence collected here (frame counts, color mix,
right-turn distribution under US rules) points to A+B. A further difference is deliberate: the results
here are over three maneuvers, without the U-turn, which follows the left signal in 99.9\% of the labels.
For these reasons the scores of this thesis and the published ones are close in scale but are not a
head-to-head comparison.

\paragraph{Structured pipelines.} The original paper reports that its direction-aware model clearly
outperforms structured pipelines. Its strongest structured baseline, TLD-READY~\cite{polley2024tldready},
was designed for per-light relevance on German roads and adapted to the four-maneuver output with a
fallback to the circular signal; it reaches a TSS of 0.842 and a False-Go rate of 0.302. The structured
model of this thesis reaches 0.899 and 0.152 on subset A val, because its classes encode the maneuvers
each light controls and therefore mirror the label rules exactly, so that a perfect detector would be a
perfect predictor. Allowing for the differences in labels and directions, this suggests that the gap
reported in the paper reflects the design of the structured baselines it used rather than an inherent
limit of the structured approach. This is in line with the broader experience in detection that
well-tuned, strongly pretrained detectors~\cite{zhao2024detrs,chen2024lwdetr,robinson2025rfdetr} are hard
to beat.

\section{Theoretical and practical implications}
\label{sec:disc-implications}

\paragraph{Aggregate scores hide the safety-critical cases.} The right-turn analysis shows that a model
can reach state-of-the-art TSS while failing on nearly every forbidden right turn, because these cases
are rare and pooled metrics weight them by frequency. The same holds, to a lesser degree, for yellow and
for straight \emph{none}. Benchmarks for this task should therefore report per-direction results,
class-wise recall of the rare states and, ideally, a False-Go rate per maneuver. The TSS, which weighs
False-Go and False-Stop equally and lets the trivial right direction contribute a quarter of the
accuracies, is a poor proxy for safety on its own.

\paragraph{Label conventions are part of the model.} The meaning of a right-turn label differs between
the US and Germany, and the choice of convention changes the right-turn distribution from 97\% \emph{none}
to 63\% green. An end-to-end model internalizes whichever convention it is trained on. A structured
pipeline can change the rule at inference time, which is a practical argument for keeping the traffic
rules explicit, especially for systems meant to operate in more than one country.

\paragraph{Detection is the bottleneck out of distribution.} The ATLAS results show that most errors of a
transferred model begin with lights that are not detected, and that detection depends steeply on light
size. In the structured pipeline, a missed red light becomes a \emph{go} decision by default, which
explains the high False-Go on Argoverse~2 images. A safer design would treat ``no relevant light
found'' near an intersection as uncertain rather than as permission, as the map-based ATLAS framework
does by mapping an unknown signal-group state to red~\cite{polley2025atlas}. The same work shows that a
tele camera and a wide-angle camera are the practical answer to the size problem, and that a detector
trained on DTLD alone is clearly weaker on ATLAS images than one that also saw ATLAS data (mAP@50 about
0.53 versus 0.7); the per-light results obtained
here with a different detector point the same way. Trinci et al.~\cite{trinci2026color} likewise found
that relevance models do not transfer between DTLD and US data, so the transfer gaps measured in this
thesis appear to be a property of the data rather than of a particular architecture.

\paragraph{Cost and deployability.} The structured model runs at 22\,ms per frame on a consumer GPU and
needs only a few epochs of fine-tuning, which makes it a realistic component for an embedded perception
stack. Its explicit intermediate output (which lights were found and what they control) also supports
monitoring and debugging, an important property for safety arguments.

\section{Limitations}
\label{sec:disc-limits}

\begin{itemize}
  \item \textbf{Reconstructed labels.} The benchmark is a reconstruction. Its rules were validated against
        the published distribution and by spot-checking about 80 frames across datasets, but no full
        manual quality control was done. Some label noise certainly remains, and the left~$=$~none
        mismatch shows a systematic difference from the original.
  \item \textbf{Single seeds.} Each configuration was trained once. The epoch-to-epoch variation of the
        OpenLane-V2 model (about one TSS point, \cref{tab:rfdetr-epochs}) is of the same order as the
        differences one would want to resolve, so conclusions about small differences are limited.
  \item \textbf{Evaluation on the validation split.} OpenLane-V2's test split has no usable annotations,
        so, as in the original paper, the validation split is the evaluation set. Model selection was
        kept away from it, but it is still a split that has been looked at during development.
  \item \textbf{No head-to-head comparison with the original model.} Neither the original labels nor
        its code are public, so the published numbers can only be set beside the results of this thesis
        with the caveats of \cref{sec:disc-compare}. The structured model's class encoding is, moreover,
        tailored to the label rules.
  \item \textbf{Size-based relevance on LISA.} LISA has no relevance labels, so the ego intersection's
        lights were identified by their size. Trinci et al.~\cite{trinci2026color} show that the
        largest light is a poor relevance cue in the United States; the LISA labels are therefore
        noisier than those of OpenLane-V2 and DTLD, and the LISA results should be read as an
        approximate measure of transfer.
  \item \textbf{Convention mismatch in Singapore.} For comparability, US rules were applied to left-hand
        traffic scenes, which is semantically incorrect for those frames.
  \item \textbf{Single frames.} All models see one image. Flashing signals, occlusions and the short
        yellow phase cannot be resolved reliably without temporal context.
\end{itemize}

\section{Future research directions}
\label{sec:disc-future}

\begin{itemize}
  \item \textbf{Safety-oriented training and evaluation.} Class-balanced or cost-sensitive losses that
        penalize False-Go more than False-Stop, and benchmark protocols with per-direction and
        per-class reporting, would align optimization with the actual risk.
  \item \textbf{Rule-aware models.} Conditioning the model on the applicable traffic convention (country,
        left- or right-hand traffic), or keeping the rules explicit as in the structured pipeline, would
        let a single model operate across jurisdictions and label the Singapore scenes correctly.
  \item \textbf{Temporal models.} Aggregating evidence over a short sequence would stabilize predictions,
        handle occlusions, and allow flashing and phase-change states to be modeled.
  \item \textbf{Multi-dataset training.} The DTLD and LISA labels built here use compatible conventions,
        so mixed training across countries and illumination conditions is the natural next experiment,
        with LISA's night sequences as the main robustness test. ATLAS, BSTLD, S$^2$TLD and BDD100K
        could add pictograms, cameras, weather and countries to the detector's training data, and
        VZC-TLD, the only US dataset with relevance labels, would allow an image-level test of the
        stop/go decision outside Germany.
  \item \textbf{Small-object detection.} Higher input resolution around the horizon, tiled inference or
        a dedicated tele camera (as on ATLAS) target the steepest failure mode observed here.
  \item \textbf{Uncertainty.} Calibrated confidence for each maneuver, and an explicit ``unknown'' output
        when no relevant light is found at an approaching intersection, would allow a planner to act
        conservatively instead of defaulting to \emph{go}.
  \item \textbf{Fusion with maps and topology.} Where maps or online topology
        predictions~\cite{li2023graph,wu2024topomlp} are available, they could supply the relevance that
        single-image models must currently infer.
\end{itemize}
